% from ega3-1-fr.tex
\oldpage[III]{107}
\begin{proof}
En effet, $f$ est projectif (II, 5.5.5, (v)) et comme $Y$ est projectif
sur $k$, la condition supplémentaire de
(\hyperref[III.2.3.6-fr]{2.3.6}) est vérifiée~; il suffit alors de
considérer le sous-schéma fermé $Y'$ de $Y$ défini par l'Idéal cohérent
$\mathcal J$ du cor. (\hyperref[III.2.3.6-fr]{2.3.6}).
\end{proof}

\begin{env}[2.3.8]
\label{III.2.3.8-fr}
\emph{Remarque.} Au chap. IV, en étudiant la notion de diviseur, nous
verrons que si, dans l'énoncé de
(\hyperref[III.2.3.5-fr]{2.3.5}), on suppose que les anneaux
$\mathcal O_y$ $(y\in Y)$ sont factoriels (ce qui est le cas par exemple
si $Y$ est non singulier), alors $X$ peut se déduire de $Y$ en faisant
éclater un sous-préschéma fermé $Y'$ de $Y$ dont l'espace sous-jacent est
contenu dans $Y-U$.
\end{env}

\subsection{Une généralisation du théorème fondamental}
\label{subsection:III.2.4-fr}

\begin{theorem}[2.4.1]
\label{III.2.4.1-fr}
Soient $Y$ un préschéma noethérien, $\mathcal S$ une
$\mathcal O_Y$-Algèbre quasi-cohérente de type fini. Soient
$f:X\to Y$ un morphisme projectif, $\mathcal S'=f^*(\mathcal S)$,
$\mathcal M$ un $\mathcal S'$-Module quasi-cohérent de type fini.
Alors~:
\begin{enumerate}
\item[(i)] Pour tout $p\in\mathbb Z$, $R^pf_*(\mathcal M)$ est un
$\mathcal S$-Module de type fini.
\item[(ii)] Soit de plus $\mathcal L$ un $\mathcal O_X$-Module inversible
ample pour $f$, et posons
$\mathcal M(n)=\mathcal M\otimes\mathcal L^{\otimes n}$ pour tout
$n\in\mathbb Z$. Il existe un entier $N$ tel que, pour $n\geq N$, on ait
\[
R^pf_*(\mathcal M(n))=0
\tag{2.4.1.1}\label{III.2.4.1.1-fr}
\]
pour tout $p>0$, et que l'homomorphisme canonique
\[
f^*(f_*(\mathcal M(n)))\longrightarrow\mathcal M(n)
\]
(0\textsubscript{I}, 4.4.3) soit surjectif.
\end{enumerate}
\end{theorem}

\begin{proof}
Posons $Y'=\operatorname{Spec}(\mathcal S)$,
$X'=\operatorname{Spec}(\mathcal S')$ de sorte que
$X'=X\times_Y Y'$ (II, 1.5.5)~; soient $g:Y'\to Y$,
$g':X'\to X$ les morphismes structuraux, qui sont affines par
définition, et $f'=f_{(Y')}:X'\to Y'$~; on a donc un diagramme
commutatif
\[
\xymatrix{
X\ar[d]_f & X'\ar[l]_{g'}\ar[d]^{f'}\ar[dl]^{h}\\
Y & Y'\ar[l]^{g}
}
\]
et le morphisme $f'$ est projectif (II, 5.5.5, (iii))~; posons
$h=f\circ g'=g\circ f'$.

(i) Soit $\widetilde{\mathcal M}$ le $\mathcal O_{X'}$-Module associé au
$\mathcal S'$-Module quasi-cohérent $\mathcal M$, quand $X'$ est
considéré comme un $X$-schéma affine (II, 1.4.3), de sorte que l'on a
$\mathcal M=g'_*(\widetilde{\mathcal M})$~; comme $\mathcal M$ est un
$\mathcal S'$-Module de type fini, $\widetilde{\mathcal M}$ est un
$\mathcal O_{X'}$-Module de type fini (II, 1.4.5)~; comme $h$ est de type
fini, puisque $g$ et $f'$ le sont (II, 1.3.7 et I, 6.3.4, (ii)), $X'$ est
noethérien (I, 6.3.7) et $\widetilde{\mathcal M}$ est par suite cohérent.
Cela étant, comme $g'$ est affine, l'homomorphisme canonique
$R^pf_*(\mathcal M)\to R^ph_*(\widetilde{\mathcal M})$ est bijectif
(\hyperref[III.1.3.4-fr]{1.3.4}). En outre, cet homomorphisme est un
homomorphisme de $\mathcal S$-Modules~; en effet, de l'homomorphisme
canonique
\[
g'_*(\mathcal O_{X'})\otimes_{\mathcal O_X}
g'_*(\widetilde{\mathcal M})\longrightarrow
g'_*(\widetilde{\mathcal M})
\tag{2.4.1.2}\label{III.2.4.1.2-fr}
\]
qui définit la structure de $\mathcal S'$-Module de $\mathcal M$ (en se
rappelant que $\mathcal S'=g'_*(\mathcal O_{X'})$), on déduit
canoniquement un homomorphisme
\[
f_*(g'_*(\mathcal O_{X'}))\otimes
R^pf_*(g'_*(\widetilde{\mathcal M}))\longrightarrow
R^pf_*(g'_*(\widetilde{\mathcal M})).
\]

